\section*{الفصل \ref{ch:b3:astrophysics} --- الفيزياء الفلكية}

\begin{solution}{exo:b3:astrophysics:1}
(أ) $P_{\text{c}} \sim GM^2/R^4 \approx \qty{1.1e15}{Pa}$ --- أي عشرة
مليارات جوّ. (ب) $\bar\rho \approx \qty{1400}{kg/m^3}$: ومنه
$T \sim P m_{\text{p}}/\rho k_{\text{B}} \approx \qty{e8}{K}$ ---
وهو تقدير خشن، أعلى برتبة من القيمة الحقيقية \qty{1.6e7}{K}، لكن المقياس
صحيح. (ج) وهو عند عتبة النفق تمامًا. (د) والنجم الأوّلي
ينكمش ويسخن \emph{حتى} يشتعل اللبّ، ثم يتوقف عن
الانكماش: فكل كرة غاز ضخمة بما يكفي لا بدّ أن تبلغ الاشتعال
--- فالنجم يضبط نفسه على منظِّم الحرارة النووي.
\end{solution}

\begin{solution}{exo:b3:astrophysics:2}
(أ) $L = 4\pi d^2\times1360 \approx \qty{3.8e26}{W}$. (ب) $T =
(L/4\pi R^2\sigma)^{1/4} \approx \qty{5800}{K}$. (ج) وبفين:
$\lambda_{\text{ذروة}} \approx \qty{500}{nm}$ --- فقد تطورت العين
على ذروة الشمس. (د) وتدفق التربيع العكسي (السنة الأولى)؛ و
ستيفان--بولتزمان وفين: أي جسم \cref{ch:b3:photons-phonons}
الأسود، مستنتَجًا من إحصاء الفوتونات.
\end{solution}

\begin{solution}{exo:b3:astrophysics:3}
(أ) سيريوس B: أقصى اليسار وأقصى الأسفل؛ وبيتلجوز: أقصى اليمين وأقصى الأعلى.
(ب) $R/R_\odot = \sqrt{L/L_\odot}\,(T_\odot/T)^2$: فسيريوس B
$\approx 0.009\,R_\odot \approx$ نصف قطر الأرض؛ وبيتلجوز
$\approx 860\,R_\odot \approx \qty{4}{au}$. (ج) وعطارد والزهرة و
الأرض والمريخ كانت لتدور \emph{داخله}. (د) وبحجم الأرض ومع ذلك
بكتلة الشمس: فلا شيء إلا ضغط انحلال الإلكترونات
(\cref{ch:b3:quantum-statistics}) --- فسيريوس B هو \omterm{ex:b3:quantum-statistics:dwarf}{القزم
الأبيض} الكلاسيكي.
\end{solution}

\begin{solution}{exo:b3:astrophysics:4}
(أ) $10\,\text{مليار سنة}\times10^{-2.5} \approx \qty{30}{Myr}$. (ب)
و $0.5^{-2.5} \approx 5.7$: أي نحو 57 مليار سنة --- فكل قزم
أحمر وُلد قط لا يزال يحترق. (ج) والإضاءة تنمو وفق
$M^{3.5}$ أما الخزان فوفق $M$ لا غير: فالأغنياء يحرقون ثروتهم
أسرع مما يودعونها. (د) والنجوم الضخمة تموت في غضون بضعة
ملايين سنة من ولادتها --- قبل أن تنجرف إلى أيّ مكان: فتنفجر المستعرات العظمى
داخل الأذرع الحلزونية المكوِّنة للنجوم التي صنعتها.
\end{solution}

\begin{solution}{exo:b3:astrophysics:5}
(أ) $E = -K$: فالإشعاع يجعل $E$ أشدّ سلبيةً، ومنه ترتفع $K$ ---
أي درجة الحرارة --- \emph{ارتفاعًا}: وهو توقيع الثقالة
الترموديناميكي. (ب) $GM^2/R \approx \qty{3.8e41}{J}$. (ج)
$E/L_\odot \approx \qty{e15}{s} \approx \qty{30}{Myr}$. (د) و
شمس كلفن لم يكن ليتجاوز عمرها عشرات ملايين السنين، بينما
كانت ديدان داروين وصخور لايل تطالب بالمليارات: و
كان الحلّ وقودًا لم يعرفه أحد --- أي \cref{ch:b3:nuclear-physics}،
قبل سبعين سنة.
\end{solution}

\begin{solution}{exo:b3:astrophysics:6}
(أ) $\bar\rho \approx \qty{1.8e9}{kg/m^3}$: أي ملعقة صغيرة
بتسعة أطنان. (ب) $g = GM/R^2 \approx \qty{3.3e6}{m/s^2}$ --- أي ثلاثمئة
ألف ضعف ثقالة الأرض. (ج) و $n \propto M/R^3$:
فالانحلال يعطي $P \propto M^{5/3}/R^5$، والثقالة تطلب
$P \propto M^2/R^4$؛ وبالمساواة، $R \propto M^{-1/3}$. (د) والوزن
الأكبر يحتاج إلى ضغط فيرمي أكبر، ولا يوفّره إلا الانضغاط
--- فتُكمِش الكتلة النجمَ \emph{إكماشًا}؛ وللّولب
جدار: فما إن تصير الإلكترونات نسبوية حتى يليّن قانون الضغط إلى
$n^{4/3}$، ويخفق التوازن عند أيّ نصف قطر، وتكون
$M_{\text{Ch}} = 1.4\,M_\odot$ حافة الجرف.
\end{solution}

\begin{solution}{exo:b3:astrophysics:7}
(أ) $V = M/\rho_{\text{نووي}} \approx \qty{1.2e13}{m^3}$:
ومنه $R \approx \qty{14}{km}$ --- أي نواة لها أفق مدينة. (ب)
و $\omega \propto 1/R^2$: فالتسريع الدوراني بمقدار $(10^8/1.4\times10^4)^2
\approx 5\times10^{7}$، ومنه تنكمش مدة 25 يومًا إلى
${\sim}\qty{40}{ms}$: أي أرض النوابض. (ج) $g \approx
\qty{e12}{m/s^2}$: فجبل هنا يكون ميليمترات هناك.
(د) $B \approx 10^{-2}\times5\times10^{7} \approx
\qty{5e5}{T}$ --- والمحور المغناطيسي، وهو غير محاذٍ لمحور
الدوران، يقمع الإصدار الراديوي على طول قطبيه: فتمسح الحزمة
مسح المنارة، ونسمّي الومضات نابضًا.
\end{solution}

\begin{solution}{exo:b3:astrophysics:8}
(أ) $\tfrac12mc^2 = GMm/R$ يعطي $R = 2GM/c^2$. (ب) وللشمس:
\qty{3.0}{km}؛ وللأرض: \qty{8.9}{mm} --- فاضغط أيًّا منهما داخل
ذلك ولن يستطيع الضوء المغادرة. (ج) $R_{\text{s}} \approx
\qty{1.2e10}{m}$: أي خُمس مدار عطارد --- وتُشاهَد نجوم
تسوط حوله. (د) $\bar\rho \propto
M/R_{\text{s}}^3 \propto 1/M^2$: فثقب بكتلة $10^8\,M_\odot$ أقلّ
كثافة من الماء. فلا حاجة إلى ``سحق'' شيء عند الأفق ---
إذ هو نقطة لا عودة، لا سطح صخر.
\end{solution}

\begin{solution}{exo:b3:astrophysics:9}
(أ) $H_0 = 7\times10^4/3.09\times10^{22} \approx
\qty{2.3e-18}{s^{-1}}$. (ب) $1/H_0 \approx \qty{4.4e17}{s}
\approx 14$ مليار سنة. (ج) $v = \qty{1.4e4}{km/s}$؛ و
$z \approx v/c \approx 0.047$: فتصل الأطوال الموجية ممطوطةً
5\,\%. (د) والمعدّل لم يكن ثابتًا --- إذ أبطأته
المادة باكرًا، وسرّعته الطاقة المظلمة منذئذٍ: وكون
$1/H_0$ يقع مع ذلك في حدود مليار سنة من القيمة الحقيقية 13.8 مصادفة
كونية صغيرة.
\end{solution}

\begin{solution}{exo:b3:astrophysics:10}
(أ) $\lambda_{\text{ذروة}} = b/T \approx \qty{1.1}{mm}$:
أي ميكرويّات، كما أُعلن. (ب) وممطوطة بمقدار $3000/2.725 \approx
1100$ --- فالكون أكبر بألف ومئة مرة في كل
اتجاه منذئذٍ. وإعادة الالتحام تنتظر عند \qty{3000}{K}، أي دون
$\qty{13.6}{eV}/k_{\text{B}}$ بكثير، لأن توازن ساها
(\cref{prop:b3:grand-canonical:saha}) يميله ملياران من
الفوتونات لكل \omterm{def:b3:particle-physics:zoo}{باريون}: فالذيل المؤيِّن يبقي الهدروجين مفكَّكًا
طويلًا بعد ما تقوله الطاقات الساذجة. (ج) وهي نحو
$1.6\times10^{9}$ فوتونًا لكل \omterm{def:b3:particle-physics:zoo}{باريون} --- فالكون، بالعدّ،
ضوء كله تقريبًا. (د) والخلفية ليست مكانًا بل
السماء كلها: فكل خط رؤية ينتهي على ضباب الانتثار
الأخير، ومنه يستقبلها كل هوائي، مصوَّبًا إلى أيّ مكان.
\end{solution}

\begin{solution}{exo:b3:astrophysics:11}
(أ) $n/p = \eu^{-1.29/0.8} \approx 0.20$. (ب) ومع $n/p = 1/7$:
يكون 14 بروتونًا لكل نيوترونين؛ والنيوترونان يخطفان بروتونين إلى
$^4$He واحد (كتلته 4) فيتبقّى 12 بروتونًا وحيدًا: أي $4/16 = 25\,\%$
بالكتلة. (ج) ولم يكن لأيّ نجم وقت ليجعل ربع كل شيء
هليومًا؛ ووجود تلك الأرضية في أقدم الغازات وأفقرها بالمعادن
هو الكون الفتيّ الساخن متلبّسًا --- أي عامل بولتزمان
مكتوبًا عبر السماء. (د) ولا توجد نواة مستقرة
عند $A = 5$ أو $8$، وكانت الكثافة ودرجة الحرارة
تهبطان في غضون دقائق: فالجسر إلى الكربون (بثلاثة هليومات
دفعةً واحدة) لا ينفتح إلا في ألباب العمالقة الحمر الكثيفة الصبورة.
\end{solution}

\begin{solution}{exo:b3:astrophysics:12}
(أ) $mv^2/r = GM(r)m/r^2$. (ب) $M = v^2r/G \approx
\qty{1.1e42}{kg} \approx 5\times10^{11}\,M_\odot$. (ج) أي نحو
العُشر: فتسعون بالمئة من المجرّة لا تضيء ولا
تمتصّ. (د) و $v$ المسطَّحة تعني $M(r) \propto r$ --- فالكتلة
غير المرئية تظلّ تتراكم وراء القرص المرئي بكثير، في هالة
مظلمة ممتدة (كثافتها $\sim1/r^2$) لا تكون فيها المجرّة
المضيئة إلا اللبّ المنار.
\end{solution}

\begin{solution}{pb:b3:astrophysics:1}
\textbf{1.} $t \sim 10\,\text{مليار سنة}\times20^{-2.5} \approx
\qty{6}{Myr}$: فستعيش الشمس ألفَي ضعف عمره.
\textbf{2.} النوى الأثقل تحمل شحنة أكبر: فتحتاج حواجز كولوم
الأعلى إلى ألباب أسخن كي تجاري احتمالات النفق
(\cref{exo:b3:nuclear-physics:8}).
\textbf{3.} الحديد يجلس على قمّة منحني $B/A$: فاندماجه
\emph{يكلّف} طاقة --- وينتهي سلّم وقود الفرن على رماد.
\textbf{4.} إنها \omterm{prop:b3:astrophysics:compact}{كتلة تشاندراسيخار}: أي أقصى ما يستطيع ضغط
انحلال الإلكترونات حمله؛ فاللبّ قزم أبيض نما
داخل نجم حيّ، وعند $1.4\,M_\odot$ يخفق سنده.
\textbf{5.} الهدروجين: ملايين السنين؛ والهليوم: مئات الألوف؛ و
السيليكون: يوم --- فكل مرحلة تدفع أقلّ لكل كيلوغرام بينما
ينفق النجم، وهو يزداد سخونةً، أسرع فأسرع (وأكثره نيوترينوات):
فتجنّ الساعة جنونًا لوغاريتميًّا.
\textbf{6.} $t \sim 1/\sqrt{G\rho} =
1/\sqrt{6.7\times10^{-11}\times10^{12}} \approx \qty{0.1}{s}$:
فينسحب اللبّ من تحت النجم في عُشر ثانية.
\textbf{7.} $E \sim GM^2/R \approx \qty{4e46}{J}$.
\textbf{8.} وعرض الفوتونات هو $10^9\times3.8\times10^{26}
\times2.6\times10^{6} \approx \qty{e42}{J}$ --- أي جزء من عشرة
آلاف. والباقي يغادر نيوترينوات: فالمادة النووية
المنهارة معتمة على كل ما عداها، لكنها
(\cref{ch:b3:particle-physics}) شبه شفافة لجسيمات القوة
الضعيفة نفسها، وهي التي تصرف الطاقة في ثوانٍ.
\textbf{9.} $E/4\pi d^2 = 4\times10^{46}/2.8\times10^{43}
\approx \qty{1.4e3}{J/m^2}$.
\textbf{10.} وعند $\qty{15}{MeV} = \qty{2.4e-12}{J}$ لكلٍّ:
يكون ${\sim}6\times10^{14}$ نيوترينو لكل متر مربع --- أي عبر
كل متر مربع من الأرض.
\textbf{11.} $6\times10^{14}\times10^{-46}\times1.5\times
10^{32} \approx 9$: وقد سجّل الكاشف إحدى عشرة --- أي
إصابة مباشرة في رتبة القدر لنجم مات قبل 160\,000 سنة.
\textbf{12.} النيوترينوات تغادر اللبّ عند الانهيار؛ أما
الضوء فلا يولد إلا حين تخترق الصدمة غلاف النجم
بعد ساعات. وكون كليهما وصل في الليلة نفسها بعد
160 ألف سنة من الطيران يحدّ سرعة النيوترينو إلى سرعة الضوء
في حدود أجزاء من $10^{9}$ --- ويحدّ كتلته إلى ما لا يكاد يُذكر.
\textbf{13.} ميزانية الطاقة ($\sim10^{46}\,\unit{J}$، و 99\,\%
منها نيوترينوات) ومدة الرشقة البالغة ${\sim}10$ ثوانٍ ---
كلاهما طابق النظرية في انهيار اللبّ، فرقّى
آلية \omterm{rem:b3:astrophysics:evolution}{المستعر الأعظم} من السبورة إلى المشاهدة.
\textbf{14.} $\bar\rho = 2.8\times10^{30}/(\tfrac43\pi
(1.2\times10^4)^3) \approx \qty{4e17}{kg/m^3}$: أي الكثافة
النووية --- أي ملعقة \cref{def:b3:nuclear-physics:nuclide} الصغيرة،
وقد صارت بحجم مدينة.
\textbf{15.} التسريع الدوراني $(10^8/1.2\times10^4)^2 \approx
7\times10^{7}$: فتصير مدة 30 يومًا $\to$ ${\sim}\qty{40}{ms}$.
\textbf{16.} $B \approx 10^{-2}\times7\times10^{7} \approx
\qty{7e5}{T}$ --- أي مئة مليون تسلا-متر من التدفق،
محفوظةً في طابع بريد.
\textbf{17.} المحور المغناطيسي، وهو مائل عن محور الدوران،
يشعّ الراديو حزمةً على طول قطبيه؛ والحزمة تمسح السماء مسح
المنارة. وبقيّة السرطان --- أي نجم 1054 الضيف --- لا تزال
تدور ثلاثين مرة في الثانية، ولم يكد ألف سنة من تباطؤ الدوران يبدأ.
\textbf{18.} $R_{\text{s}} = 3\times\qty{3}{km} \approx
\qty{9}{km}$ --- أي من رتبة النجم نفسه: فمع انعدام
أيّ ضغط يجادل، يستمر الانهيار عبر الأفق
وتكون البقيّة ثقبًا أسود.
\textbf{19.} مصفوفات من ساعات كهذه كواشف لأمواج الثقالة
بفتحة بحجم المجرّة (وهي أفضل مخابر
النسبية: إذ تضمر مدارات النوابض الثنائية تمامًا كما
يتنبأ الإشعاع الثقالي).
\textbf{20.} أكسجينه اندمج في قشرة نجم ضخم
وقذفه مستعر أعظم؛ والحديد صيغ في
النهاية الانفجارية نفسها؛ وانجرف كلاهما في السحب بين النجوم، ثم
انضمّا إلى السديم الشمسي، فالأرض، فالبحر، فالسلسلة الغذائية
--- فالقارئ. أما هدروجين الانفجار العظيم فقد كان
لك منذ البدء.
\textbf{21.} كتلة انفجار واحدة $\to$ سطوع ذاتي واحد
$\to$ والسطوع الظاهري يقرأ البعد: أي شمعة
معلومة القدرة تُرى عبر مليارات السنين الضوئية ---
وهي المسطرة التي ترسم الطرف البعيد من مخطط هابل.
\textbf{22.} التمدد \emph{يتسارع}: فثمة
``طاقة مظلمة'' تدفع الكون متباعدًا --- أي نوبل 2011،
وأكبر فاتورة مفتوحة في الفيزياء.
\textbf{23.} $\qty{50}{kpc} \approx 163000$ سنة ضوئية: فقد
غادر الضوء بينما كان \emph{الإنسان العاقل} يغادر إفريقيا أول مرة.
\textbf{24.} $6\times10^{14}\times0.03 \approx 2\times10^{13}$
عبر كل إنسان حيّ --- ومع ذلك، باحتمالات
\cref{exo:b3:particle-physics:10}، لم يتفاعل شخص من كل
مليون مع واحد منها: فالومضة التي فاقت مجرّةً سطوعًا مرّت بالبشرية بلا إحساس.
\textbf{25.} نجم توازن عشرة ملايين سنة على بروتونات
تنفق؛ وانهيار في عُشر ثانية دفع $10^{46}$
جول نيوترينوات؛ ومبدأ باولي يمسك الجثة عند
الكثافة النووية؛ والحطام المبعثر --- الأكسجين والحديد و
القارئ --- لا يزال يركب تمدد هابل، في كون
صورته الوليدة منحني بلانك. فالفيزياء، مكتملة الآن:
أُغلق الكتاب.
\end{solution}
